\documentclass[12pt,a4paper]{article}
\usepackage[utf8]{inputenc}
\usepackage[indonesian]{babel}
\usepackage[margin=2.5cm]{geometry}
\usepackage{tabularx}
\usepackage{parskip}

\begin{document}

\begin{flushleft}
    {\large \textbf{KEPOLISIAN NEGARA REPUBLIK INDONESIA}} \\
    {\large \textbf{DAERAH TIRTA JAYA }} \\
    {\large \textbf{RESOR KOTA TIRTA JAYA}}
    \rule{\textwidth}{1.5pt}
    \vspace{0.3cm}
    {\textbf {"PRO JUSTITIA"}}
\end{flushleft}
    
\begin{center}
    {\large \textbf{BERITA ACARA PEMERIKSAAN}}\\
    {\large \textbf{TERSANGKA}} \\
    \textbf{Nomor: BAP/34/II/2026/Reskrim}
    \vspace{0.3cm}
\end{center}

---

Pada hari ini, Jumat, tanggal 22 September 2026, pukul 14.00 WIB, bertempat di Ruang Unit Reserse Kriminal Polres Metro Kota Tirta Jaya, saya:
\begin{tabbing}
    \hspace{3cm} \= \hspace{1cm} \= \kill
    \textbf{Nama Penyidik} \> : \> Liya \\
    \textbf{Pangkat/NRP} \> : \> Inspektur Satu Polisi NRP 77012345 \\
    \textbf{Jabatan} \> : \> Penyidik Reserse Kriminal Polres Metro Kota Tirta Jaya
\end{tabbing}

Berdasarkan Surat Perintah Penyidikan Nomor: \textbf{SP.Sidik/45/II/2026/Reskrim}, telah melakukan pemeriksaan terhadap seorang yang diduga sebagai Tersangka:
\begin{tabbing}
    \hspace{3cm} \= \hspace{1cm} \= \kill
    \textbf{Nama Lengkap} \> : \> Hendra Wijaya \\
    \textbf{Tempat/Tgl. Lahir} \> : \> Jakarta, 12 Mei 1969 \\
    \textbf{Umur} \> : \> 57 Tahun \\
    \textbf{Pekerjaan} \> : \> PNS (Kepala Dinas PUPR Kota Tirta Jaya / Pengguna Anggaran) \\
    \textbf{Alamat} \> : \> Jl. Mawar No. 10, Kota Tirta Jaya \\
    \textbf{Agama} \> : \> Islam \\
    \textbf{Kebangsaan} \> : \> Indonesia
\end{tabbing}

Selanjutnya disebut sebagai \textbf{TERSANGKA}.

\vspace{0.3cm}
\begin{center}
    \textbf{--- HASIL PEMERIKSAAN ---}
\end{center}

\textbf{Pertanyaan 1} \\
\textbf{Penyidik :} Apakah Saudara dalam keadaan sehat jasmani dan rohani saat ini, serta bersedia memberikan keterangan dengan sebenar-benarnya terkait perkara ini? \\
\textbf{Tersangka :} Ya, saya dalam keadaan sehat dan bersedia memberikan keterangan dengan sebenar-benarnya.

\vspace{0.2cm}
\textbf{Pertanyaan 2} \\
\textbf{Penyidik :} Apakah Saudara mengetahui alasan Saudara dipanggil dan diperiksa oleh penyidik pada hari ini? \\
\textbf{Tersangka :} Ya, saya tahu. Saya diperiksa terkait keterlibatan saya sebagai Kepala Dinas PUPR dan Pengguna Anggaran (PA) dalam proyek pembangunan Sistem Pengolahan Air Bersih (SPAB) Kota Tirta Jaya tahun anggaran 2021.

\vspace{0.2cm}
\textbf{Pertanyaan 3} \\
\textbf{Penyidik :} Bagaimana proses penunjukan Rizal Pratama sebagai Pejabat Pembuat Komitmen (PPK) serta penunjukan PT Delta Konsultan Infrastruktur dalam proyek ini? \\
\textbf{Tersangka :} Penunjukan Saudara Rizal Pratama sebagai PPK sudah melalui pertimbangan administratif di dinas, sedangkan penunjukan PT Delta Konsultan Infrastruktur dilakukan untuk menyusun dokumen perencanaan teknis proyek.

\vspace{0.2cm}
\textbf{Pertanyaan 4} \\
\textbf{Penyidik :} Berdasarkan penyelidikan, ditemukan adanya pertemuan tertutup di sebuah hotel di Jakarta antara Saudara, Rizal Pratama, dan Arman Gunawan (Direktur PT Sumber Tirta Konstruksi) sebelum tender dimulai. Apa yang dibahas dalam pertemuan tersebut? \\
\textbf{Tersangka :} Pertemuan tersebut hanya koordinasi biasa terkait kesiapan proyek, tidak ada pembahasan mengenai pengaturan pemenang tender.

\vspace{0.2cm}
\textbf{Pertanyaan 5} \\
\textbf{Penyidik :} Bagaimana tanggapan Saudara terhadap hasil audit Badan Pemeriksa Keuangan (BPK) yang menyatakan bahwa proyek ini mengalami kerugian negara sebesar Rp34.000.000.000 akibat pengurangan spesifikasi dan pekerjaan fiktif? \\
\textbf{Tersangka :} Saya hanya bertanggung jawab secara administratif sebagai Pengguna Anggaran. Saya tidak terlibat dalam teknis pelaksanaan di lapangan maupun pengawasan pekerjaan tersebut.

\vspace{0.2cm}
\textbf{Pertanyaan 6} \\
\textbf{Penyidik :} Apakah ada pihak lain yang turut memberikan instruksi atau menekan Saudara dalam proses pengadaan dan pelelangan proyek SPAB ini? \\
\textbf{Tersangka :} Tidak ada pihak yang menekan saya secara langsung. Semua proses berjalan sesuai dengan alur birokrasi di dinas.

\vspace{0.2cm}
\textbf{Pertanyaan 7} \\
\textbf{Penyidik :} Apakah seluruh keterangan yang Saudara berikan di atas sudah benar dan tidak ada paksaan dari pihak manapun? \\
\textbf{Tersangka :} Ya, semua keterangan tersebut sudah benar.

\vspace{0.5cm}
\begin{center}
    \textbf{--- PENUTUP ---}
\end{center}

Demikian Berita Acara Pemeriksaan ini dibuat dengan sebenar-benarnya dan setelah dibacakan kembali kepada tersangka, ia menyatakan telah memahami isi keterangannya, serta menandatangani BAP ini tanpa tekanan atau paksaan dari pihak manapun.

\vspace{0.5cm}

\begin{tabularx}{\textwidth}{X c}
    & Dibuat di : Tirta Jaya \\
    & Pada Tanggal : 22 September 2026 \\
    & \\
    \textbf{Tersangka,} & \textbf{Penyidik,} \\
    \vspace{1.5cm} & \\
    \underline{\textbf{Hendra Wijaya}} & \underline{\textbf{Liya}} \\
    & Inspektur Satu Polisi NRP 77012345
\end{tabularx}

\end{document}
